% PA194FW mathematical proof fragment.  Requires only amsmath, amssymb,
% amsthm and booktabs; all additional notation is local to this fragment.
% The stack-compiler/QF-induction organization is adapted, with attribution,
% from references/AX11H/ax11h_preservation.tex.  The fused axiom, ten-entry
% interface, explicit-binder generalization and associated certificates are new.
\section{PA194FW: a fused negative-implication axiom}
\label{sec:pa194fw}
\providecommand{\FWI}{\mathsf I}
\providecommand{\FWU}{\mathsf U}
\providecommand{\FWM}{\mathsf M}
\providecommand{\FWE}{\mathsf E}
\providecommand{\FWN}{\mathsf N}
\providecommand{\FWC}{\mathcal C}
\providecommand{\FWPRA}{\mathrm{iPRA}}
\newtheorem{fwtheorem}{Theorem}[section]
\newtheorem{fwlemma}[fwtheorem]{Lemma}

PA194FW has 194 working states, two tape symbols, one tape and a separate
halt state.  Its source changes the axiom interface itself: two previously
separate axiom heads become one negative-implication head.  This reduces
the number of entries produced by AX from eleven to ten.  A new grammar
and ten checked state identifications retain a total of 194 states.  The
search did not produce a 193-state machine and provides no lower bound.

There are two distinct ancestors.  The implementation starts with the
symmetrically cached controller of PA194H5XS, whose fifth head is
$\FWI(\FWU(D,\FWN(\FWE(H,D))),A)$.  The preservation argument
below compares directly with the already calibrated ordinary calculus
$\FWC_4$ of PA194AX11H, whose fifth head is
$\FWN(\FWU(D,\FWN(\FWE(H,D))))$.  A replacement of one arbitrary
source word by another would not establish that comparison; both the
formula translation and the exact ten-entry source interface are proved.

\subsection{Numerical formulas and the exact change}

Use the inherited total numerical constructors
\[
 p(a,b)=\frac{(a+b)(a+b+1)}2+b,\qquad
 \FWI(a,b)=4p(a,b),\qquad \FWM(a,b)=4p(a,b)+1,
\]
\[
 \FWU(i,a)=4p(a,i)+2,\qquad \FWE(a,b)=4p(a,b)+3,
 \qquad \FWN(a)=\FWI(a,1).
\]
Write $a\to b$ for $\FWI(a,b)$, with right-associated arrows.  Code 1
is contradiction.  The numerical alias $\FWI(0,0)=0$ is retained.
Define $C(a,b)=\FWN(\FWI(a,\FWN(b)))$ and
$\exists_i a=\FWN(\FWU(i,\FWN(a)))$.  The ordinary heads are
\begin{align*}
 h_1(D,H)&=\FWI(\exists_D\FWU(D,H),H),\\
 h_2(D,A,H)&=\FWI(\FWU(D,\FWU(A,H)),\FWU(A,\FWU(D,H))),\\
 h_3(D,A,H)&=\FWI(\FWM(H,D),\FWI(\FWE(H,A),\FWM(A,D))),\\
 h_4(D,A,H)&=\FWI(\FWE(H,D),\FWI(\FWM(A,H),\FWM(A,D))),\\
 V(D,H)&=\FWU(D,\FWN(\FWE(H,D))),\qquad h_5(D,H)=\FWN(V(D,H)),\\
 h_6(D,A,H)&=\FWI(\FWU(D,A),
                   \FWI(\FWU(D,\FWI(A,H)),\FWU(D,H))).
\end{align*}
The new head is the literal numerical term
\[
 \boxed{F_{45}(D,A,H)=
       \FWN\bigl(\FWI(h_4(D,A,H),V(D,H))\bigr).}       \tag{FW-F}
\]
It is classically equivalent to the conjunction of $h_4$ and $h_5$.
It is \emph{not} literally $C(h_4,h_5)$: that abbreviation would put
$\FWN(\FWN(V))$, rather than $V$, in the second argument of the inner
implication.  All source and proof certificates use (FW-F) literally.

Let $\FWC_F$ replace the two families $h_4,h_5$ in $\FWC_4$ by the one
family $F_{45}$.  It retains $h_1,h_2,h_3,h_6$ and
\[
 \Gamma(A)=\FWI(A,\FWU(A,A)),\qquad
 \rho(D,A,H,u)=((H\to u)\to D)\to((D\to H)\to(A\to H)),
 \quad u\le D.
\]
It also retains the guarded gate, defined without any semantic shorthand by
\begin{align*}
 P(H,D,A)&=\FWI(A,\FWM(A,D)),\\
 V_g(H,D,A)&=\FWI(\FWM(H,D),\FWU(A,\FWM(H,D))),\\
 S_g(H,D,A)&=\FWU\bigl(D,C(\FWI(\FWM(H,D),P(H,D,A)),
                              \FWI(\FWM(A,D),A))\bigr),\\
 Q_g(H,D,A)&=\FWU(A,\FWI(S_g(H,D,A),\FWN(P(H,D,A)))),\\
 G(H,D,A)&=\FWN(\FWI(V_g(H,D,A),Q_g(H,D,A))),
                  \qquad A\ne H,\ A\ne D.
\end{align*}
The inference rules are modus ponens and theorem generalization
$a\mapsto\FWU(i,a)$.  A failed gate guard emits the provable code zero;
indeed $\rho(0,0,0,0)=0$.  These are numerical proof systems over all
natural codes, not a free syntactic datatype with a forbidden zero alias.

\subsection{Noncircular fusion and projection proofs}

The implicational schemas $K$, $S$ and Peirce used in the following
certificates are derived rules, not new primitive axioms.  The pinned file
\texttt{references/logic/RestrictedImplication.lean} contains explicit
finite derivations \texttt{k}, \texttt{s} and \texttt{peirce}.  Its
\texttt{Rules} interface uses only MP, the restricted $\rho$ row,
reflexivity of the parameter bound, and the two implications
\[
 u\le a\ \Longrightarrow\ u\le\FWI(a,b),\qquad
 u\le b\ \Longrightarrow\ u\le\FWI(a,b).
\]
The numerical constructor dominates both arguments, so these side
conditions hold in $\FWPRA$.  Neither $h_4$, $h_5$, $F_{45}$ nor any
quantifier axiom occurs in that interface.  Substituting natural codes
in these fixed derivations gives primitive recursive proof constructors;
checking their finitely many equations and bound implications uses
quantifier-free induction only.  Thus using $K,S,$ and Peirce below
does not assume the desired calibration.

The retained first head is also essential to the particular proof given
here.  For arbitrary $a$, put $Y=\FWU(0,\FWN(\FWU(0,a)))$.
Then $h_1(0,a)$ is $(Y\to1)\to a$.  The $K$ instance
$1\to(Y\to1)$, followed by that head and the usual $K,S$ deduction
transformation, proves
\[
                         1\to a.                         \tag{FW-Ex}
\]
No extra explosion axiom is silently inserted.  In particular, replacing
the ordinary first head by the joined first head from the J-family is
not justified by this argument.

\begin{fwlemma}[Fixed syntactic certificates]
In both calculi the following formulas have uniform finite derivations:
\[
 a\to\bigl(\FWN(b)\to\FWN(\FWI(a,b))\bigr),\qquad
 \FWN(\FWI(a,b))\to a,\qquad
 \FWN(\FWI(a,b))\to\FWN(b).                              \tag{FW-Cert}
\]
Their proof-list constructors and checkers are primitive recursive and
their validity is provable in $\FWPRA$.
\end{fwlemma}
\begin{proof}
For the first formula, under assumptions $a$, $b\to1$, and $a\to b$,
two applications of MP yield 1; discharge the three assumptions using
$K,S$.  For the third, assumptions $b$ and $(a\to b)\to1$ yield 1
using $K$ to obtain $a\to b$; discharge both assumptions.  For the
second, assume $(a\to b)\to1$.  Under the additional assumption
$a\to b$, MP followed by (FW-Ex) gives $a$.  Discharge that additional
assumption, apply Peirce $((a\to b)\to a)\to a$, and discharge the
original assumption.  These arguments are expanded by the delivered
checker into proof lists of respectively 161, 205 and 53 lines.  It
checks every axiom substitution, every earlier-premise index and every
MP equality; it uses no truth-table oracle.  Expansion of the derived
$K,S,$ Peirce schemas into restricted-$\rho$ proofs gives proofs in the
stated primitive calculus.  Each expansion is a fixed finite expression,
so uniform numerical substitution and proof concatenation have the
claimed primitive recursive and quantifier-free verification properties.
\end{proof}

\begin{fwlemma}[Equivalence of the numerical proof systems]
There are primitive recursive proof transformations in both directions
between $\FWC_4$ and $\FWC_F$, preserving the conclusion.  $\FWPRA$
proves their correctness for every internally checked finite proof.
\end{fwlemma}
\begin{proof}
Translate a $\FWC_4$ leaf $h_4(D,A,H)$ by the left projection of
$F_{45}(D,A,H)$.  Translate $h_5(D,H)$ by the right projection of
$F_{45}(D,0,H)$.  Thus sharing the three parameters within the fused
head causes no loss of instances: every desired old fourth-head instance
has its own left projection, and every fifth-head instance has a right
projection with, for example, $A=0$.  In the opposite direction, use
the old $h_4(D,A,H)$ and $h_5(D,H)$ in the first certificate of
(FW-Cert).  All other leaves and both inference rules are retained.

For a proof represented as a finite sequence with earlier premise
indices, scan its lines and append the indicated finite replacement
proof, recording the output index for each original line.  Replace MP
and generalization indices by these recorded indices.  This is a
primitive recursion on the input length.  Validity of the translated
prefix and its index map is a bounded, hence primitive recursive Boolean
predicate.  Quantifier-free induction proves the invariant; numerical
substitution in the fixed certificates supplies the leaf cases.  This
proves preservation for arbitrary natural proof codes and yields, in
$\FWPRA$,
\[
 \exists e\,\operatorname{Prf}_{\FWC_4}(e,1)
 \quad\longleftrightarrow\quad
 \exists f\,\operatorname{Prf}_{\FWC_F}(f,1).             \tag{FW-Logic}
\]
\end{proof}

\subsection{The exact ten-entry command interface}

Write $[a;w]=p(a,w)$ and extend this notation by repeated pairing with
the first displayed item outermost.  The projections are total, with
$\operatorname{hd}(0)=\operatorname{tl}(0)=0$ and $[0;0]=0$.
Let the store be $(t,w,d,h)$.  With $A=c$, $D=d+x$, $H=h+y$, the new
AX command is
\[
 \mathsf A_{c,x,y,z}(t,w,d,h)
       =(0,[q_0,q_1,\ldots,q_8,t;w],z,0),               \tag{FW-AX}
\]
where the exact top-to-bottom entries are
\[
\begin{array}{c|l}
 j&q_j\\\hline
 0&G(H,D,A),\text{ or zero if }A=H\text{ or }A=D\\
 1&\Gamma(A)\\
 2&h_6(D,A,H)\\
 3&F_{45}(D,A,H)\\
 4&h_3(D,A,H)\\
 5&h_2(D,A,H)\\
 6&h_1(D,H)\\
 7&\rho(D,A,H,d)\\
 8&\FWU(A,h).
\end{array}
\]
The ninth generated entry has the explicit command binder $A=c$ and
the \emph{incoming} cache $h$.  Substituting $\FWU(d,h)$ here would
give the older AX11H interface and would be incorrect for this source.
All $q_0,\ldots,q_7$ have residue zero modulo four; $q_8$ has residue
two and is at least two.  No generated entry equals 1.

The other commands, unchanged from the cached-test interpreter, are
\[
 \mathsf S(t,w,d,h)=
 \begin{cases}\mathsf{halt},&h=1,\\
 (h,\operatorname{tl}w,d,\operatorname{hd}w),&h\ne1,
 \end{cases}                                             \tag{FW-S}
\]
and, for $b=d+x$,
\[
 \mathsf{MP}_x(t,w,d,h)=
 \begin{cases}(b,w,b,0),&\FWI(t,b)=h,\\
 (0,w,b,0),&\FWI(t,b)\ne h.
 \end{cases}                                             \tag{FW-MP}
\]
Only $\mathsf S$ halts.  The fourth AX field is consumed by its
fallthrough through the complete MP reader: after the gate has been
pushed, $t=d=h=0$, so comparing $\FWI(0,z)$ with zero leaves $t=h=0$
on both branches and sets $d=z$.  Thus (FW-AX) holds for every fourth
field, not just the zero fields used below.

\subsection{Stack words and their early-halting cases}

Composition is from left to right.  Put $A_0=\mathsf A_{0,0,0,0}$ and
define
\begin{align*}
 \mathsf{Push}&=A_0;\mathsf S^9;\mathsf{MP}_0,&
 \mathsf{Pop}&=A_0;\mathsf S^{11},\\
 \mathsf{Discard}&=\mathsf{MP}_0,&
 \mathsf{Offset}_i&=\mathsf{MP}_i.
\end{align*}
The exact contracts are
\begin{align*}
 \mathsf{Push}(t,w,d,h)&=(0,[t;w],0,0),\\
 \mathsf{Pop}(t,[a;w],d,h)&=
   \begin{cases}\mathsf{halt},&t=1,\\(t,w,0,a),&t\ne1,\end{cases}\\
 \mathsf{Discard}(t,w,0,h)&=(0,w,0,0),\\
 \mathsf{Offset}_i(0,w,0,0)&=(0,w,i,0).                 \tag{FW-Words}
\end{align*}
For example, after $A_0;\mathsf S^9$ the top is $q_7$, the cache is
$q_8$, and the stack is $[t;w]$; the test for successful MP cannot
confuse a residue-zero implication with the residue-two cache.  In any
case the consequent is zero, so Push has the displayed result.  After
$A_0;\mathsf S^{10}$ the top is $q_8$ and the cache is the saved $t$;
the next Shift tests precisely that saved value.  These observations
also cover $t=1$ and the zero stack alias.

For $0\le j\le7$, let
\[
 \mathsf{Sel}_j(A,D,H)=\mathsf{Discard};
 \mathsf A_{A,D,H,0};\mathsf S^{j+2};\mathsf{Pop}^{8-j}.
                                                               \tag{FW-Sel}
\]
From $(t,w,0,0)$ it returns $(q_j,w,0,0)$, with the indicated
parameters and incoming offset zero.  After the initial shifts the
selected head is top; repeated Pop discards the remaining generated
entries and finally the saved zero.  The selected head is not 1, so
those Pops cannot halt.  For the restricted row $\rho(D,A,H,u)$,
$u\le D$, use instead
\[
 \mathsf{Discard};\mathsf{Offset}_u;
 \mathsf A_{A,D\mathbin{\dot-}u,H,0};\mathsf S^9;\mathsf{Pop}.
                                                               \tag{FW-Rho}
\]
The truncated subtraction is legitimate because
$u+(D\mathbin{\dot-}u)=D$ under the displayed bound.

The explicit binder in $q_8$ gives a shorter generalization recipe:
\[
 \mathsf{Gen}_i=\mathsf{Push};\mathsf S;
             \mathsf A_{i,0,0,0};\mathsf S^{10},
 \qquad
 \mathsf{Gen}_i(a,w,0,0)=(\FWU(i,a),w,0,0).              \tag{FW-Gen}
\]
This also holds when $a=1$: the only time 1 is cached is immediately
before AX clears the cache, so it is not tested by Shift at that point.

The MP suffix is
\[
 \mathsf S;\mathsf S;\mathsf{MP}_b;\mathsf{Pop}.          \tag{FW-MPS}
\]
On $(0,[a,\FWI(a,b),0;w],0,0)$ it halts if $a=1$ or $b=1$,
and otherwise returns $(b,w,0,0)$.  The two Shifts expose $a$ and
$\FWI(a,b)$; MP obtains $b$; Pop removes the reserved zero and tests
the saved conclusion.  The reserved zero is necessary.  All these
word identities are checked symbolically in
\texttt{verification/check\_command\_algebra.py}; its 19 contracts
include the unit-minor and unit-conclusion cases.  The checker uses total
Cantor projection identities, not an injective free-list representation.

\subsection{A uniform proof-to-word compiler in iPRA}

A numerical proof is a finite sequence with earlier premise indices.
Compile its lines to words $W_j$ by primitive recursion.  Use (FW-Sel)
and (FW-Rho) for axiom leaves.  Append (FW-Gen) to the premise word for
generalization.  For MP with major $e_1$ proving $\FWI(a,b)$ and minor
$e_0$ proving $a$, use
\[
 W_e=\mathsf{Discard};\mathsf{Push};W_{e_1};\mathsf{Push};
 W_{e_0};\mathsf{Push};\mathsf S;\mathsf S;
 \mathsf{MP}_b;\mathsf{Pop}.                             \tag{FW-Compile}
\]
The invariant is
\[
 W_e(t,w,0,0)=\mathsf{halt}\quad\hbox{or}\quad
 W_e(t,w,0,0)=(\operatorname{concl}(e),w,0,0).            \tag{FW-Cmp}
\]
The initial Discard and Push reserve a zero; the two premise computations
and their Pushes produce exactly the precondition of (FW-MPS).  An earlier
halt remains a valid outcome of every extension.  For a refutation $e$,
\[
                          A_0;W_e;\mathsf{Pop}             \tag{FW-Ref}
\]
therefore halts from every incoming store.  The first AX normalizes the
offset and cache; the final Pop tests the returned contradiction.  A
terminal plain Shift would test zero and would not suffice.

The arbitrary-base-stack quantifier in (FW-Cmp) must not be smuggled into
an induction formula of iPRA.  Here is a bounded verification.  For a
proof of $N>0$ lines let $L$ bound all compiled word lengths and $K$ bound
their command codes.  For a total coded command step $\operatorname{next}$
with an absorbing halt flag, put
\[
 \Phi_K(s)=1+\max\left(s,
       \max_{c\le K,\ v\le s}\operatorname{next}(c,v)\right).
\]
It is a monotone primitive recursive majorant.  Quantifier-free induction
proves that every prefix of at most $L$ commands, started below $s$,
remains below $\Phi_K^L(s)$.  For an arbitrary normalized input code
$s_0$ define $B(j)=\Phi_K^{L(N-1-j)}(s_0)$.  If $i<j$ then
$\Phi_K^L(B(j))\le B(i)$.  Induct on the Boolean conjunction asserting
(FW-Cmp) for every previous line $i<j$ on every normalized store code
at most $B(i)$.  This is a primitive recursive Boolean test, since each
word evaluation and every quantifier is bounded.  In the MP and
generalization cases, each call to an earlier word follows a prefix of
at most $L$ commands, so its actual input lies in the corresponding
earlier bounded domain.  The local contracts prove the new line.  At
the last line $B(N-1)=s_0$; generalizing the arbitrary parameter $s_0$
proves (FW-Cmp) without strengthening the induction schema.  This is the
bounded-input method of the attributed AX11H proof, with every word and
entry count changed to the explicit ten-entry interface above.

\subsection{Exhaustive scheduling and syntactic soundness}

Let $e(\epsilon)=0$ and $e(c::w)=1+b_c(e(w))$, where
\begin{align*}
 b_{\mathsf S}(z)&=p(0,z),\\
 b_{\mathsf{MP}_x}(z)&=p(1,p(x,z)),\\
 b_{\mathsf A_{c,x,y,u}}(z)&=p(c+2,p(y,p(x,p(u,z)))).      \tag{FW-Code}
\end{align*}
The exact field order is $y,x,u$.  A live scheduler state is
$(n,\ell,t,w,d,h)$, starting at $(2,0,0,0,0,0)$.  It selects seed $n$
and advances the fresh counter to $n+1$ if $\ell=0$; otherwise it selects
$\ell-1$ and leaves the fresh counter alone.  It executes one decoded
command and sets the continuation to the residual after all command
fields.  Stores are carried between candidates.

Every residual is a nested Cantor component and is at most the raw seed.
Positive continuations strictly decrease, so a fresh candidate $N$ drains
in at most $N+1$ command invocations, unless it halts first.  Given the
all-store halting suffix $V_e$ of (FW-Ref), prefix another $A_0$.
Its fresh raw seed is
\[
 N_e=p\bigl(2,p(0,p(0,p(0,e(V_e))))\bigr)\ge3.
\]
This candidate executes the prefix and then the exact suffix; the prefix
may alter the store because the suffix works from every store.  The
absorbing scheduler has halted by $K(e)=(N_e+2)^2$ commands.  The
decoding identities and draining bound are primitive recursive arithmetic
statements with quantifier-free induction proofs.  No reset of a failed
candidate and no least-successful-candidate search is assumed.

For soundness maintain a finite valid $\FWC_F$ derivation list and
indices proving the top, cache and canonical stack entries.  For
$w>0$, $\operatorname{tl}(w)<w$, so decoding the canonical stack takes
at most $w$ iterations.  The zero stack supplies virtual zero entries
using $\rho(0,0,0,0)$.  AX appends proofs of its eight axiom heads, of
$\FWU(A,h)$ by generalization of the incoming cache proof, and of the
saved top.  Its restricted row satisfies $d\le d+x$; a failed gate
uses the fixed zero proof.  MP combines the existing top and cache
proofs when the comparison succeeds and otherwise returns zero.  A
nonhalting Shift transports existing proof indices.  A halting Shift
exposes a proof of the cache value 1.  All updates and validity tests are
primitive recursive; a Boolean quantifier-free induction on the command
count extracts a refutation from every scheduler halt.  Consequently
\[
 \FWPRA\vdash\operatorname{Halt}(R_F)
       \longleftrightarrow
       \exists e\,\operatorname{Prf}_{\FWC_F}(e,1).      \tag{FW-Scheduler}
\]

\subsection{Binding the source and its clocks}

The grammar has 191 macro tokens.  In the H5XS emitter the end of the
old fifth-head block was \texttt{copyA;cons}, which constructs
$\FWI(V(D,H),A)$.  Replace it by
\texttt{cons;push;incT;cons}.  At that point the top is $V(D,H)$
and the saved stack head is $h_4(D,A,H)$, so the first cons makes
$\FWI(h_4,V)$; push followed by incT and cons negates it.  This
consumes the separate fourth head and leaves exactly (FW-F).

This local explanation is accompanied by a complete free-term check of
the entire source word.  The checker verifies all intermediate scratch
preconditions, every final register, both failed-gate cases and the
unconsumed fourth field $p(z,\mathrm{tail})$ before the MP tail.  In
particular it verifies the explicit-binder entry $\FWU(A,h)$.  An
independent arithmetic-template construction matches all 3,695 primitive
instruction rows, arguments and both outgoing branches, and a separately
resolved named binary source tree matches the same graph.  Its ports are
\[
\begin{array}{c|rrrrrrrr}
 &\mathrm{fetch}&\mathrm{decode}&\mathrm{AX}&\mathrm{MP}&
 \mathrm{compare}&\mathrm{success}&\mathrm{return}&\mathrm{Shift}\\\hline
 \mathrm{site}&25&34&50&3599&3669&3679&3686&3687.
\end{array}
\]
The result binds (FW-AX)--(FW-MP) to the delivered source, not merely to
a proposed abstract interpreter.

For completeness, the source clock is an actual primitive recursive
function.  The retained arithmetic templates have exact durations
\[
\begin{array}{c|c}
 \mathrm{template}&\mathrm{duration}\\\hline
 \mathrm{increment}&1\\
 \mathrm{clear}(r)&2r+1\\
 \mathrm{transfer}(r,s)&3r+1\\
 \mathrm{duplicate}(r,s,t)&4r+1\\
 \mathrm{copy}(r)&7r+2\\
 \mathrm{times4}(y)&6y+1\\
 \mathrm{pair}(t,a)&4t+2+7(a+t)(a+t-1)/2+5(a+t).
\end{array}
\]
The subtraction in the polynomial is interpreted at zero by its natural
value zero.  Define $(a_0,b_0)=(0,0)$ and
$(a_{j+1},b_{j+1})=(a_j-1,b_j+1)$ when $a_j>0$, and
$(b_j+1,0)$ otherwise.  For unpairing input $n$, the exact duration is
\[
 2+3b_n+\sum_{j<n}
       \begin{cases}5,&a_j>0,\\3b_j+8,&a_j=0.\end{cases}
\]
The comparison duration is $3h+2$ for $h=x$, $3x+2$ for $h<x$,
and $2h+x+4$ for $h>x$.  Fetch costs 2 for a continuation and
$7n+4$ for a fresh seed.  Compose these functions in literal source
order, updating their register arguments after each template.  This is
exactly \texttt{Arithmetic} and \texttt{command\_clock} in the source
checker.  Summing the resulting positive duration along the coded
scheduler recursion defines its cumulative clock $B(k)$.

The uniform proofs are ordinary counter inductions: transfer and
duplicate decrease the source counter; the pairing loop uses the
triangular sum; the unpairing recurrence proves its pair index equals
$j$ and sums the displayed step costs; comparison decreases both
counters until one is exhausted.  Frame preservation is part of each
template's register equation.  The complete proofs and precise loop
equations are retained with attribution in
\texttt{references/AX11H/source\_command\_completion.tex}.
They apply unchanged to these templates and to any finite macro word;
only the checked literal composition is new.  The 882 local numerical
tests and 324 composed Shift/MP tests are additional regression evidence,
not replacements for these uniform induction arguments.

\subsection{The literal machine and physical transport}

Construction first builds 135 fixed states and 69 reader states.  It
then performs one fixed-controller identification and nine simultaneous
reader identifications, with all 40 local predecessor-marker contexts
checked on the relevant final stage of the table.  Thus the delivered
table has $135+69-1-9=194$ working states.  The complete literal
verifier is independent of the construction compiler and arithmetic-box
oracle.  It reconstructs four generic unary/list sweep lemmas from the
literal rows, then replays every one of the 4,809 source outcomes over
nine arbitrary natural registers, arbitrary affine nonnegative gap and
an arbitrary nonempty finite spare list.  It checks the complete tape
frame, including both blank exteriors.  The replay takes 1,130,294
symbolic steps; the largest contract takes 403 such steps.

Every nonhalting outcome has zero net frame motion and the independently
resolved successor address after padding to 48 bits.  There is no
address overflow; the longest source path has length 38.  Every affine
duration has nonnegative coefficients and constant at least 3.  The
unique halt instruction has duration 3.  These are exhaustive finite
bindings to the generic literal-row induction lemmas, not a finite
sample of register values or a tape-window simulation.

From the blank tape, native state 74 and head position zero, the machine
reaches source entry 25 in exactly 57,768 physical steps, in native
state 84 with head at $-2$, logical registers $(2,0,\ldots,0)$, and 26
unit spares.  The startup checker compares the whole finite support of
the tape with the exact source-frame encoding.  Dense renaming sends
the initial state to 0 and the separate halt state to 194.  All 194
working states are reachable in the directed transition graph; this is
a graph claim, not a claim that every state is visited from the blank
tape under the actual run.

To internalize the physical argument, encode the full finite tape
support, head, control, frame gap and spare list by natural numbers.
Each of the four sweep lemmas has a counter/list induction proof with
a primitive recursive Boolean frame invariant.  Instantiate them in
the finite replayed contracts and compose by quantifier-free induction
on source steps.  The resulting cumulative physical clock $\tau(j)$
includes the changing gap, list width and spare values; they are not
held fixed.  Its step increment is zero at the canonical absorbing halt and positive
at every live source step.  The supplied
\texttt{references/AX11H/physical\_ipra\_audit.tex} and
\texttt{ipra\_simulation\_theorem.tex} give these general frame and
positive-clock transport proofs.  Their finite premises are precisely
the source and physical certificates just checked for PA194FW.

For example, a numerical refutation $e$ has the explicit physical
halting bound
\[
                 57768+\tau\bigl(B((N_e+2)^2)\bigr).      \tag{FW-Time}
\]
Conversely, if the physical run halts at time $t$, positive live
durations ensure that the simulated source reaches its absorbing halt
within at most $t+1$ source steps after initialization; otherwise its
positive cumulative clock would outrun $t$ while still being live.
The same argument applies to the positive source-command clocks.
Decidability of the coded halt flag handles these finite alternatives
intuitionistically.  This proves in $\FWPRA$
\[
 \operatorname{Halt}_{\rm blank}(\mathrm{PA194FW})
 \longleftrightarrow\operatorname{Halt}(R_F).           \tag{FW-Physical}
\]

\subsection{Composition with the inherited PA calibration}

The inherited theorem is the \emph{written} AX11H calibration in
\texttt{references/AX11H/consistency\_equivalence.tex}, with its
Foundation adapter, finite-demand logical compiler, model-uniform seed
soundness and two one-way intuitionistic transfer endpoints.  It
concerns the same fixed Foundation PA-refutation predicate as the pinned
repository.  It supplies
\[
 \FWPRA\vdash
 \operatorname{Con}(\mathrm{PA})\longleftrightarrow
 \neg\operatorname{Halt}_{\rm blank}(\mathrm{PA194AX11H}).
                                                               \tag{FW-Cal}
\]
The inherited abstract and physical proofs also identify AX11H halting
with $\exists e\,\operatorname{Prf}_{\FWC_4}(e,1)$ in iPRA.
Combining (FW-Physical), (FW-Scheduler), (FW-Logic) and that execution
equivalence gives a positive halting equivalence between the two
machines, entirely through their numerical proof systems.

\begin{fwtheorem}[PA194FW preservation]
With the literal binary transition table and blank-start convention
specified here, the inherited written PA calibration and the new
proofs above yield
\[
 \boxed{\FWPRA\vdash
    \neg\operatorname{Halt}_{\rm blank}(\mathrm{PA194FW})
        \longleftrightarrow\operatorname{Con}(\mathrm{PA}).}
\]
\end{fwtheorem}
\begin{proof}
Intuitionistic logic transports negation across the just-proved halting
equivalence; compose with (FW-Cal).  No PA-consistency assumption,
Markov principle or inference of an existential from its double
negation is used.  In particular this composition does not claim a
new positive iPRA equivalence between numerical refutation existence
and Foundation PA-refutation existence.  The inherited calibration
uses its stated model-uniform and negative-transfer route.
\end{proof}

The new claim is a written mathematical preservation proof with
explicit primitive recursive constructions and fully replayed finite
machine bindings.  It is relative to the identified inherited written
absolute calibration in the ordinary sense that a new mathematical
result uses an earlier theorem.  The available Lean infrastructure and
the supplied execution checkers do not amount to an end-to-end Lean
proof of this machine's PA-consistency equivalence.  No such completed
kernel proof is claimed.  The source files, exact hashes and recorded
verification outcomes make the remaining formalization target concrete.
